% META: {"id":"1SPE-EXPO-CO-047","chapitre":"1SPE-EXPONENTIELLE","type_objet":"corrige","exercice_id":"1SPE-EXPO-EX-047","status":"generated"}
% BEGIN-VERIFY
% from sympy import *
% x = symbols('x', real=True)
% f = exp(2*x) - 2*exp(x)
% fp = diff(f, x)
% assert simplify(fp - 2*exp(x)*(exp(x) - 1)) == 0
% assert solve(fp, x) == [0]
% assert f.subs(x, 0) == -1
% END-VERIFY
\begin{corrige}{1SPE-EXPO-EX-047}
\begin{enumerate}
\item On dérive $f(x) = \mathrm{e}^{2x} - 2\mathrm{e}^x$ :
\[f'(x) = 2\mathrm{e}^{2x} - 2\mathrm{e}^x = 2\mathrm{e}^x(\mathrm{e}^x - 1).\]

\item Le facteur $2\mathrm{e}^x > 0$ pour tout réel $x$.

Le signe de $f'(x)$ est donc celui de $\mathrm{e}^x - 1$ :
\begin{itemize}
\item $\mathrm{e}^x - 1 < 0 \iff \mathrm{e}^x < 1 \iff x < 0$ ;
\item $\mathrm{e}^x - 1 = 0 \iff x = 0$ ;
\item $\mathrm{e}^x - 1 > 0 \iff x > 0$.
\end{itemize}

Donc $f$ est strictement décroissante sur $]-\infty\,;0]$ et strictement croissante sur $[0\,;+\infty[$.

\item On calcule :
\[f(0) = \mathrm{e}^0 - 2\mathrm{e}^0 = 1 - 2 = -1.\]

D'après le tableau de variations, $f$ admet un \textbf{minimum} égal à $\boxed{-1}$, atteint en $x = 0$.

\item $f(-1)=\mathrm{e}^{-2}-2\mathrm{e}^{-1}\approx-0{,}60$ et $f(1)=\mathrm{e}^{2}-2\mathrm{e}\approx1{,}95$. Avec le minimum $(0;-1)$ et les variations précédentes, ces points permettent d'esquisser la courbe.
\end{enumerate}
\end{corrige}
